% Small, maintainable exercise include file.
% Responsibilities are split into includes under `exercises.d/`:
% - `input-path.tex`      : sets `\input@path` for subvariant inputs
% - `setup-counter.tex`   : ensures the `exercise` counter exists
% - `include-exercise.tex`: defines `\IncludeExercise{<path>}` wrapper
% The actual exercise inclusions are then simple calls to \IncludeExercise{<path>}.

% Load path, counter and the include wrapper (keeps this file minimal)
% Ensure TeX looks for subvariant files in the exercise folder when main.tex uses relative \input
\makeatletter
\def\input@path{{../../../ExerciseDatabase/matematica/P4_funcoes/4-funcao_inversa/determinacao_analitica/MAT_P4FUNCOE_4FX_DAX_003/}}
\makeatother

% Define exercise counter if not already defined (defensive)
\makeatletter
\@ifundefined{c@exercise}{\newcounter{exercise}}{}
\makeatother

% Define a simple wrapper command to include an exercise file and
% print a standardized heading. The wrapper toggles the style flag
% `\showexerciciotitle` (defined in style.tex) so the included
% exercise does not duplicate its own heading.
\providecommand{\IncludeExercise}[1]{%
  \begingroup
    \refstepcounter{exercise}%
    \noindent\textbf{Exercício \theexercise.}\par
    \showexerciciotitlefalse
    \input{#1}%
    \showexerciciotitletrue
  \endgroup
}


% Exercicio 1: Complementa as equações, de forma a que exista igualdade entre os números de diferentes sistemas numéricos, i.e. (10)_10 = (20)_5.

\exercicio{
    Complementa as equações, de forma a que exista igualdade entre os números de diferentes sistemas numéricos, i.e. $(10)_{10} = (20)_5$, pois $(10)_{10} = 1*10^1 + 0*10^0 = 10$ e $(20)_5 = 2*5^1 + 0*5^0 = 10$.
\begin{enumerate}[label=(\alph*)]
    \item $(15)_{10} = (  \hspace{3cm} )_2$
    \item $(100)_{10} = (  \hspace{3cm} )_8$
    \item $(255)_{10} = ( \hspace{3cm}  )_{16}$
    \item $(1111)_2 = ( \hspace{3cm}  )_{10}$
    \item $(777)_8 = ( \hspace{3cm}  )_{10}$
    \item $(FF)_{16} = ( \hspace{3cm}  )_{10}$
\end{enumerate}
}

\exercicio{
Considere a sucessão dos números triangulares representados visualmente abaixo
(os valores $T_1,T_2,T_3,T_4$ estão indicados nas figuras):
\begin{center}
\begin{tikzpicture}[scale=0.5]
% T1 = 1
\foreach \x in {0} {
    \fill (\x,0) circle (0.15);
}
% Valor $T_1 = 1$ representado na figura

% T2 = 3
\begin{scope}[xshift=3cm]
\foreach \x in {0,1} {
    \fill (\x,0) circle (0.15);
}
\fill (0.5,0.5) circle (0.15);
% Valor $T_2 = 3$ representado na figura
\end{scope}

% T3 = 6
\begin{scope}[xshift=6cm]
\foreach \x in {0,1,2} {
    \fill (\x,0) circle (0.15);
}
\foreach \x in {0.5,1.5} {
    \fill (\x,0.5) circle (0.15);
}
\fill (1,1) circle (0.15);
% Valor $T_3 = 6$ representado na figura
\end{scope}

% T4 = 10
\begin{scope}[xshift=10cm]
\foreach \x in {0,1,2,3} {
    \fill (\x,0) circle (0.15);
}
\foreach \x in {0.5,1.5,2.5} {
    \fill (\x,0.5) circle (0.15);
}
\foreach \x in {1,2} {
    \fill (\x,1) circle (0.15);
}
\fill (1.5,1.5) circle (0.15);
% Valor $T_4 = 10$ representado na figura
\end{scope}
\end{tikzpicture}
\end{center}
\begin{enumerate}[label=(\alph*)]
    \item Represente os quatro primeiros termos da sucessão $(T_1,T_2,T_3,T_4)$ em sistema binário.
    \vspace{3cm}
    \item Escreva o 5.º termo da sucessão $(T_5)$ em notação decimal e hexadecimal.
    \vspace{3cm}
    \item Determine o 10.º termo da sucessão $(T_{10})$ e represente-o em sistema binário.
    \vspace{3cm}
\end{enumerate}
}

\exercicio{
Observe a sucessão dos números quadrados representados visualmente abaixo
(os valores $Q_1,Q_2,Q_3,Q_4$ estão indicados nas figuras):
\begin{center}
\begin{tikzpicture}[scale=0.4]
% Q1 = 1
\foreach \x in {0} {
    \foreach \y in {0} {
        \fill (\x,\y) circle (0.12);
    }
}
% Valor $Q_1 = 1$ representado na figura

% Q2 = 4
\begin{scope}[xshift=3cm]
\foreach \x in {0,1} {
    \foreach \y in {0,1} {
        \fill (\x,\y) circle (0.12);
    }
}
% Valor $Q_2 = 4$ representado na figura
\end{scope}

% Q3 = 9
\begin{scope}[xshift=7cm]
\foreach \x in {0,1,2} {
    \foreach \y in {0,1,2} {
        \fill (\x,\y) circle (0.12);
    }
}
% Valor $Q_3 = 9$ representado na figura
\end{scope}

% Q4 = 16
\begin{scope}[xshift=12cm]
\foreach \x in {0,1,2,3} {
    \foreach \y in {0,1,2,3} {
        \fill (\x,\y) circle (0.12);
    }
}
% Valor $Q_4 = 16$ representado na figura
\end{scope}
\end{tikzpicture}
\end{center}
\begin{enumerate}[label=(\alph*)]
    \item Represente os quatro primeiros termos da sucessão $(Q_1,Q_2,Q_3,Q_4)$ em sistema hexadecimal.
    \vspace{3cm}
    \item Escreva o 5.º termo da sucessão $(Q_5)$ em notação decimal e binária.
    \vspace{3cm}
    \item Refira a definição recursiva da sucessão dos números quadrados.
    \vspace{3cm}
    \item Escreva a expressão geral (forma fechada) para o n-ésimo termo da sucessão dos números quadrados.
    \vspace{3cm}
\end{enumerate}
}